\section{Dopamine Dependent Insight Threshold}
To provide a biological basis for the developed Latent Hazard Model, we integrated the COMT Val158Met polymorphism into the model. The COMT enzyme is responsible for dopamine degradation in the prefrontal cortex (PFC). Neurobiologically, dopamine is known to regulate synaptic plasticity 
through modulation of long-term potentiation (LTP) thresholds in 
prefrontal and hippocampal circuits. Higher tonic dopamine levels 
facilitate synaptic potentiation and reduce the evidence required to 
stabilize a new behavioral strategy.

Within the latent hazard framework, this mechanism can be interpreted 
as a shift in the internal evidence threshold required to trigger the 
exploration–exploitation transition. Consequently, genotypes associated 
with higher dopamine levels should reach the insight transition earlier, 
while genotypes with lower dopamine tone require more accumulated 
reward evidence before switching strategies.The literature shows that the Met allele leads to lower enzyme activity and consequently higher tonic dopamine levels \cite{bilder2004}. 
Neurobiologically, high dopamine levels lower the synaptic potentiation (LTP) threshold, thus increasing learning flexibility. This is expressed in the model by a genotype-dependent scaling factor ($f$):
\begin{equation}
\theta_{bio} = \theta \times f
\end{equation}

\begin{figure*}[t]
    \centering
\includegraphics[width=\textwidth]{A.png}
    \caption{Genotype-dependent differences in insight timing predicted by the latent hazard model. (A) Distribution of simulated insight steps across COMT Val158Met genotypes. Points represent the median insight time across multiple simulation seeds, and error bars indicate the interquartile range (IQR). Higher dopamine levels (Met/Met) correspond to a lower effective switching threshold $\theta$, leading to earlier insight transitions. Conversely, lower dopamine levels (Val/Val) delay the onset of insight. (B) Median cumulative reward trajectories for each genotype. Curves represent the median trajectory across seeds for the latent hazard model. Vertical dashed lines indicate the estimated insight transition point for each genotype. Increasing the biological threshold $\theta_{bio} = \theta \times f$ shifts the exploration–exploitation transition to later steps, producing systematically delayed learning in the low-dopamine condition.}
    \label{fig:genotype_results}
\end{figure*}

In the simulations, genotype effects were implemented by scaling the 
switching threshold parameter $\theta$ using the genotype-dependent 
factor $f$. This modification changes the amount of accumulated reward 
evidence required before the exploration–exploitation transition occurs, 
thereby shifting the timing of the insight event.
This formulation allows dopamine-dependent genetic variation to directly 
modulate the evidence threshold for insight. In this framework, lower 
values of $\theta_{bio}$ increase the probability of earlier switching 
according to the switching rule defined in
according to the switching rule defined in Equation~\ref{eq:switch}. 

Substituting the biological threshold $\theta_{bio}$ shifts the sigmoid 
switching function toward lower reward counts, effectively predicting 
earlier exploration–exploitation transitions for high-dopamine genotypes.

Thus, higher dopamine levels effectively shift the sigmoid switching 
curve toward earlier reward counts.



\begin{table}[htbp]
\centering
\caption{Genotype-dependent Scaling Factors}
\label{tab:genotypes}
\begin{tabular}{lcc}
\toprule
\textbf{Genotype} & \textbf{Dopamine Level} & \textbf{Scaling Factor (f)} \\
\midrule
Met/Met & High    & 0.6 \\
Val/Met & Normal  & 1.0 \\
Val/Val & Low     & 1.4 \\
\bottomrule
\end{tabular}
\end{table}


Simulation results (n = 10 seeds; median [Q1, Q3]) show that dopamine-dependent shifts in the switching threshold strongly influence insight timing. The high-dopamine genotype (Met/Met) reaches insight at approximately step 1972 with the lowest effective threshold $(\theta_{bio}=8.4)$, producing the earliest transition to exploitation. In contrast, the low-dopamine genotype (Val/Val), associated with higher COMT activity and reduced dopaminergic tone, reaches insight much later at approximately step 3904 $(\theta_{bio}=19.6)$. This genotype-dependent separation is clearly visible in Figure 3 as a systematic shift of the exploration–exploitation transition toward later steps.
The logarithmic slope ratio $(log_SR)$ analysis shows that all three genotypes increase the reward collection rate after insight. The fact that the $Low_DA$ group had the highest $log_SR$ value (0.960) may indicate that the behavioral effect of the insight, even if delayed, is quite strong.